\section{Conclusion}
We have presented a general framework for the range verification of neural networks with nonlinear and learnable activation functions by treating the choice of piecewise affine abstraction as an optimization problem. We bound network-wide error as a weighted linear combination of local errors and cast the problem of allocating pieces across units as a multi-choice knapsack problem. A key property of the framework is that the abstraction is computed \emph{once per network} and reused across any number of verification queries. This amortization is precisely what makes the upfront cost of optimized allocation worthwhile in practice, and it matches the structure of real certification pipelines, which issue many queries per model~\cite{SinghEtAl2019AbstractDomain}.

Our experiments on KANs spanning $60$ to $1,010$ parameters demonstrate the value of this design. \textcolor{red}{Optimized allocation produces output bounds up to $22\times$ tighter than uniform allocation, with the largest gains at small input perturbations where verification is most operationally valuable. On the MILP stage, our approach achieves up to $15.1\times$ faster solve times than uniform allocation on the largest benchmark, and this speedup is recovered after just two verification queries on the same network.} \ck{The allocation itself is guided by the surrogate bound of Theorem~\ref{thm:kan-overall-error} rather than the true propagated error, and as we discuss in Section~\ref{sec:optimal-pwa-for-entire-kan}, this bound can be loose. We still see empirical gains because the knapsack only needs the bound's relative ranking of units to be roughly right, not its absolute value to be tight.} Future work will focus on reducing abstraction overhead for smaller networks, for instance via lazy or incremental DP, \ck{validating the surrogate directly against a sampling-based estimate of the true error on our benchmarks,} and incorporating \ck{property- and} MILP complexity estimates directly into the allocation objective to further improve end-to-end solver performance.